% GENERATED FILE. Edit canonical lessons, then run npm run generate:books.
% canonical-source-hash: fnv1a64:a322678d68f57338
% canonical-lessons: TA-C131-kili, TA-C131-pura, TA-C131-antai, TA-C131-mayil, TA-C131-erumpu

\chapter{Parrot, Pigeon, Owl, Peacock, Ant}
\label{ch:ta-parrot-pigeon-owl-peacock-ant}
\hlchaptermodality{131}

\begin{chapteropening}
\textbf{By the end of this chapter:} \emph{I can say a parrot, a pigeon, an owl, a peacock and an ant.}

Complete the last lesson of chapter 131: i can say a parrot, a pigeon, an owl, a peacock and an ant.
\end{chapteropening}

\section[kiḷi]{\ta{கிளி} (kiḷi) — a parrot}

\label{lesson:TA-C131-kili}



\begin{quote}
Before the new one: say the Tamil for a frog, then the Tamil for a tortoise.
\end{quote}

\subsection*{You'll want to know: \ta{கிளி}}
\textbf{\ta{கிளி}} — \emph{kiḷi} — \textquotedblleft{}a parrot\textquotedblright{}.

\begin{grammarlens}[title={What you've built}]
One more word for this chapter.
\end{grammarlens}

\subsection*{Your turn}
\begin{itemize}
\raggedright
  \item \emph{Say it:} \emph{kiḷi}
  \item \emph{Say it:} \emph{kiḷi}, once more
  \item \emph{Say it:} say \emph{kiḷi}
  \item \emph{Recall:} say the Tamil for a frog, then the Tamil for a tortoise, then say \emph{kiḷi} again
\end{itemize}

\begin{tcolorbox}[breakable,colback=teal!4,colframe=teal!35!black,title={Before you move on}]
What does \ta{கிளி} mean? (\textquotedblleft{}A parrot\textquotedblright{}.) Say it once more.
\end{tcolorbox}
\section[puṟā]{\ta{புறா} (puṟā) — a pigeon}

\label{lesson:TA-C131-pura}



\begin{quote}
Before the new one: say the Tamil for a tortoise, then the Tamil for a parrot.
\end{quote}

\subsection*{You'll want to know: \ta{புறா}}
\textbf{\ta{புறா}} — \emph{puṟā} — \textquotedblleft{}a pigeon\textquotedblright{}.

\begin{grammarlens}[title={What you've built}]
One more word for this chapter.
\end{grammarlens}

\subsection*{Your turn}
\begin{itemize}
\raggedright
  \item \emph{Say it:} \emph{puṟā}
  \item \emph{Say it:} \emph{puṟā}, once more
  \item \emph{Say it:} say \emph{puṟā}
  \item \emph{Recall:} say the Tamil for a tortoise, then the Tamil for a parrot, then say \emph{puṟā} again
\end{itemize}

\begin{tcolorbox}[breakable,colback=teal!4,colframe=teal!35!black,title={Before you move on}]
What does \ta{புறா} mean? (\textquotedblleft{}A pigeon\textquotedblright{}.) Say it once more.
\end{tcolorbox}
\section[āntai]{\ta{ஆந்தை} (āntai) — an owl}

\label{lesson:TA-C131-antai}



\begin{quote}
Before the new one: say the Tamil for a parrot, then the Tamil for a pigeon.

Greet the day through, in Tamil's own words: \textbf{\ta{காலை} \ta{வணக்கம்}}, \textbf{\ta{மாலை} \ta{வணக்கம்}}, \textbf{\ta{இனிய} \ta{இரவு}}.
\end{quote}

\subsection*{You'll want to know: \ta{ஆந்தை}}
\textbf{\ta{ஆந்தை}} — \emph{āntai} — \textquotedblleft{}an owl\textquotedblright{}.

\begin{grammarlens}[title={What you've built}]
One more word for this chapter.
\end{grammarlens}

\subsection*{Your turn}
\begin{itemize}
\raggedright
  \item \emph{Say it:} \emph{āntai}
  \item \emph{Say it:} \emph{āntai}, once more
  \item \emph{Say it:} say \emph{āntai}
  \item \emph{Recall:} say the Tamil for a parrot, then the Tamil for a pigeon, then say \emph{āntai} again
\end{itemize}

\begin{tcolorbox}[breakable,colback=teal!4,colframe=teal!35!black,title={Before you move on}]
What does \ta{ஆந்தை} mean? (\textquotedblleft{}An owl\textquotedblright{}.) Say it once more.
\end{tcolorbox}
\section[mayil]{\ta{மயில்} (mayil) — a peacock}

\label{lesson:TA-C131-mayil}



\begin{quote}
Before the new one: say the Tamil for a pigeon, then the Tamil for an owl.
\end{quote}

\subsection*{You'll want to know: \ta{மயில்}}
\textbf{\ta{மயில்}} — \emph{mayil} — \textquotedblleft{}a peacock\textquotedblright{}.

\begin{grammarlens}[title={What you've built}]
One more word for this chapter.
\end{grammarlens}

\subsection*{Your turn}
\begin{itemize}
\raggedright
  \item \emph{Say it:} \emph{mayil}
  \item \emph{Say it:} \emph{mayil}, once more
  \item \emph{Say it:} say \emph{mayil}
  \item \emph{Recall:} say the Tamil for a pigeon, then the Tamil for an owl, then say \emph{mayil} again
\end{itemize}

\begin{tcolorbox}[breakable,colback=teal!4,colframe=teal!35!black,title={Before you move on}]
What does \ta{மயில்} mean? (\textquotedblleft{}A peacock\textquotedblright{}.) Say it once more.
\end{tcolorbox}
\section[eṟumpu]{\ta{எறும்பு} (eṟumpu) — an ant}

\label{lesson:TA-C131-erumpu}



\begin{quote}
Before the new one: say the Tamil for an owl, then the Tamil for a peacock.
\end{quote}

\subsection*{You'll want to know: \ta{எறும்பு}}
\textbf{\ta{எறும்பு}} — \emph{eṟumpu} — \textquotedblleft{}an ant\textquotedblright{}.

\begin{grammarlens}[title={What you've built}]
One more word for this chapter.
\end{grammarlens}

\subsection*{Your turn}
\begin{itemize}
\raggedright
  \item \emph{Say it:} \emph{eṟumpu}
  \item \emph{Say it:} \emph{eṟumpu}, once more
  \item \emph{Say it:} say \emph{eṟumpu}
  \item \emph{Recall:} say the Tamil for an owl, then the Tamil for a peacock, then say \emph{eṟumpu} again
\end{itemize}

\begin{tcolorbox}[breakable,colback=teal!4,colframe=teal!35!black,title={Before you move on}]
What does \ta{எறும்பு} mean? (\textquotedblleft{}An ant\textquotedblright{}.) Say it once more.
\end{tcolorbox}
